\chapter{Transport, invariance and sensitivity across bank environments}
\label{r3:transport}
An estimate that is credible in one market may be misleading in another,
even when the same statistical method is applied correctly.
Product eligibility, customer composition, competitive responses and
measurement systems can all change. The relevant question is which
parts of a causal model remain invariant, and which target-domain
information can replace the parts that do not.

\section{Selection diagrams encode changes of mechanism}
\label{r3:selection-diagram}
Let $P$ describe a source environment and $P^*$ a target environment.
A selection diagram adds a node $S_j\to V_j$ wherever the mechanism
or exogenous distribution generating $V_j$ may differ. These nodes
mark permitted differences; they are not necessarily record-inclusion
indicators or measured causes in a pooled dataset. The absence of
$S_j\to V_j$ is a substantive invariance assertion. Transportability
asks whether the target causal law is uniquely determined by the
available source observational/experimental laws, the target
observational law and this shared model class
\citep[Definitions 4--5]{pearltransport2014}.

Suppose the only population difference is in a pretreatment covariate
$X$, and $X$ separates $Y$ from the selection nodes in the graph with
incoming arrows to $A$ removed. The conditional intervention response
is then invariant by the first rule of do-calculus. Averaging it
over the target population gives
\begin{equation}
 p^*(y\mid\doop(a))
     =\int p(y\mid\doop(a),x)\,dP^*(x).
 \label{r3:transport-formula}
\end{equation}
This is the pretreatment $S$-admissibility result in
\citet[Theorem 2 and Corollary 1]{pearltransport2014}. The right
side needs a source experiment or a justified source identification
route, and source support at target covariates. The intervention and
outcome must also have the same meaning in both places.

For a project calculation in currency units, let the causal offer
effect be $\tau(x)=10+20x$ for binary $X$. If $P(X=1)=.2$ in the
pilot and $P^*(X=1)=.8$ in the intended rollout, the average effects
are $14$ and $26$. Reusing the pilot average confuses composition
with causal response. The transport formula recovers $26$ by using
the invariant conditional effects and the new composition.

Reweighting every variable with a different distribution is not
justified. Suppose instead $U$ is a stable fair binary effect modifier
with $\tau(U)=10+20U$, and the observed proxy $Z$ equals $U$ at the
source. In the target, let $P^*(Z=1\mid U=1)=1$ and
$P^*(Z=1\mid U=0)=.6$. The population effect remains $20$ because
$P^*(U)=P(U)$. Yet weighting source proxy-specific effects $(10,30)$
by the target proxy prevalence $.8$ gives $26$. The relevant diagram
has $S\to Z\leftarrow U\to Y$; conditioning on $Z$ connects the
environment to the outcome mechanism through the latent modifier.
This finite example reconstructs the distinction between a population
shift and a proxy-measurement shift in
\citet[Section 3, Examples 1--2]{pearltransport2014}.

If the changing variable is a mediator, its target distribution must
generally be taken under the intervention, $p^*(m\mid\doop(a))$,
not its observed margin $p^*(m)$. If the response mechanism itself
changes, no reweighting of source covariates alone removes that change.
General transport may combine several source experiments and target
kernels. The source's do-calculus reduction explains how such a
formula can be justified; the simple admissibility criterion is not
a complete transport algorithm for every diagram.

\section{A transported mean has two sources of statistical error}
\label{r3:transport-gradient}
For a concrete sampling model, observe independent draws from a pooled
population. Let $S=1$ denote a source trial and $S=0$ a target cohort
whose $X$ is observed. Outcomes and treatments are available in the
source. Set $q=P(S=0)>0$, $s(x)=P(S=1\mid X=x)$,
$e_a(x)=P(A=a\mid X=x,S=1)$, and
$m_a(x)=\E[Y\mid A=a,X=x,S=1]$. Assume randomized or conditionally
exchangeable source treatment, consistency, conditional mean
transportability $\E[Y(a)\mid X,S=0]=\E[Y(a)\mid X,S=1]$,
and positive $s,e_a$ on target support. Then
\[
 \theta_a=\E[Y(a)\mid S=0]=\E[m_a(X)\mid S=0].
\]
The following is a project score derivation for this specified
observed-data model:
\begin{equation}
 \phi_a(O)=\frac{1-S}{q}\{m_a(X)-\theta_a\}
 +\frac{S}{q}\frac{1-s(X)}{s(X)}
          \frac{\ind\{A=a\}}{e_a(X)}\{Y-m_a(X)\}.
 \label{r3:transport-if}
\end{equation}
Differentiating the conditional target covariate distribution gives
the first term. A perturbation of the source outcome conditional
mean must be integrated against the target covariate distribution;
Bayes' rule converts its density ratio into $(1-s)/s$ and yields
the second. The treatment-score component is zero. Both terms are
mean zero, and their inner products with arbitrary factorized scores
recover the functional derivative. With finite variance, this
verifies the nonparametric canonical gradient for the pooled model.

Independent target and source samples with fixed sizes are usually
more natural operationally. If $r(x)=dP^*_X/dP_X$ and the source
treatment is randomized, their leading variance for the transported
mean becomes
\[
 \frac{\Var_{P^*}\{m_a(X)\}}{n_*}
 +\frac{\E_P[r(X)^2\sigma_a^2(X)/e_a(X)]}{n_s},
 \quad \sigma_a^2(x)=\Var(Y\mid X=x,A=a,S=1).
\]
The denominator $n_s$ denotes the source size. This follows by
separate influence expansions for the target empirical average and
source conditional-mean estimation; independence removes the cross
covariance. A huge target warehouse does not compensate for little
source experimentation in high-weight target strata. Estimated
selection scores, source regressions and cross-fitting need their
own remainder control; the formula alone does not establish a rate.

\section{Invariant prediction can constrain causal structure}
\label{r3:invariance}
Transport starts from specified invariant mechanisms. Invariant causal
prediction uses observed environments to test candidate restrictions
on one such mechanism. Its linear model assumes that some set $S^*$
and coefficient vector $\beta^*$ satisfy, in every environment $e$,
\[
 Y^e=\mu+(X^e_{S^*})'\beta^*+\varepsilon^e,\quad
 \varepsilon^e\perp X^e_{S^*},\quad
 \varepsilon^e\sim F_\varepsilon.
\]
The coefficient, intercept and noise law are shared. In a suitable
acyclic structural model, interventions on predictors that leave
the outcome assignment and its relevant noise independence intact
make its parent set one valid $S^*$
\citep[Assumption 1 and Proposition 1]{petersicp2016}.
An arbitrary split of calendar time into ``environments'' does not
establish these conditions; a recession may change the outcome
mechanism and measurement as well as the predictors.

For each candidate subset $S$, test the null that one such common
linear mechanism exists. Let $\widehat{\mathcal S}$ be the subsets
not rejected by valid level-$\alpha$ tests and define
$\widehat S=\bigcap_{S\in\widehat{\mathcal S}}S$, taking an empty
result if none is accepted. If the candidate family contains $S^*$,
then
\begin{equation}
 P(\widehat S\subseteq S^*)\ge1-\alpha.
 \label{r3:icp-guarantee}
\end{equation}
Indeed, whenever the true subset is not rejected, the intersection
is contained in that subset. A false positive thus requires rejecting
its one true null, an event of probability at most $\alpha$. No
union bound over all tested subsets is needed for this particular
intersection guarantee. This reproduces the proof of
\citet[Theorem 1]{petersicp2016}.

The direction of inclusion is essential: this is control of false
causal inclusions, not a promise to find every parent. Weakly differing
environments may leave the answer empty. Nor is it a probability
distribution over graphs or proof that every accepted regression is
causal. Testing all subsets is exponential in the candidate dimension;
screening that removes the true set can invalidate the guarantee.
Dependence, latent confounding, nonlinear responses and interventions
on $Y$ require the corresponding extended model and a valid test,
rather than reusing the linear theorem unchanged. In a banking
application, environments should have documented intervention or
mechanism-change meanings and temporally appropriate testing units.

\section{Sensitivity quantifies a specified failure of transport}
\label{r3:transport-sensitivity}
Define the departure from mean transportability by
\[
 b_a(x)=\E[Y(a)\mid X=x,S=0]-m_a(x).
\]
Then the true target effect equals the transported effect plus
$\E[b_1(X)-b_0(X)\mid S=0]$. If substantive analysis specifies
$|b_1(x)-b_0(x)|\le\Gamma(x)$, the identification uncertainty is
at most $\E[\Gamma(X)\mid S=0]$ in either direction. The bound
has financial units and can be calibrated against observed effect
variation across comparable markets. It is an assumption-indexed
sensitivity interval, not a confidence interval for an empirically
identified parameter. Sampling error can be added using a jointly
valid procedure conditional on the chosen sensitivity restriction.

Support failure is another, separate uncertainty. For a bounded
potential outcome $Y(a)\in[y_{\min},y_{\max}]$, let $\mathcal O$
be target covariates covered by the source, with unsupported target
mass $\omega$. If the transported conditional mean is valid on
$\mathcal O$, then
\[
 \theta_a\in\left[
  \int_{\mathcal O}m_a(x)dP^*(x)+\omega y_{\min},\quad
  \int_{\mathcal O}m_a(x)dP^*(x)+\omega y_{\max}\right].
\]
Flexible extrapolation chooses one completion inside or even outside
these unrestricted bounds; it does not remove the missing support.
Domain balancing that discards uncovered customers changes the
population being answered and should report its mass and composition.

Finally, transport to a future market combines causal invariance
with uncertainty about the future environment. A forecast of future
covariate distributions is not an observed target law. Large-scale
adoption may itself change competitor prices, customer choices and
funding costs, so that the environment depends on the policy. Such
equilibrium effects need a structural market model or a design with
variation at the relevant market scale. Historical invariance tests
and account-level causal estimates can inform that model but cannot
identify the resulting equilibrium merely by being pooled together.
